\chapter{Méthodologie de l'épreuve}

\noindent\textit{Bien connaître son cours ne suffit pas : réussir l'épreuve de
Mathématiques du BEPC demande aussi de comprendre comment elle est construite, ce que
chaque mot de consigne exige, et comment gérer son temps. C'est l'objet de ce chapitre.}
\vspace{8pt}

\section{Structure de l'épreuve}

\begin{propriete}[Ce qu'il faut savoir]
L'épreuve de Mathématiques du BEPC dure \textbf{2 heures}, coefficient \textbf{4}, notée
sur \textbf{20 points}. Elle comporte toujours deux parties \textbf{indépendantes et
obligatoires} :
\begin{itemize}
\item \textbf{Partie A — Évaluation des ressources} (10 à 12,5 points) : des exercices
courts et cloisonnés (« Activités numériques » et « Activités géométriques »), qui testent
directement les connaissances et savoir-faire du cours, chapitre par chapitre.
\item \textbf{Partie B — Évaluation des compétences} (7,5 à 10 points) : une \textbf{situation
concrète unique}, avec plusieurs « tâches » à résoudre, qui demande de \textbf{mobiliser
plusieurs notions ensemble} pour répondre à un problème de la vie courante.
\end{itemize}
Une partie des points de la Partie B (souvent $0{,}75$ à $1$ point) récompense la
\textbf{présentation générale} de la copie : propreté, clarté, unités, rédaction.
\end{propriete}

\begin{remarque}
Les deux parties sont \textbf{indépendantes} : un échec sur une question de la Partie A
n'empêche pas de réussir la Partie B, et inversement. Ne jamais abandonner une partie à
cause d'une difficulté rencontrée dans l'autre.
\end{remarque}

\section{Lexique des consignes}

\begin{propriete}[Ce que chaque mot exige vraiment]
Le barème sanctionne autant la \textbf{qualité de la réponse au mot employé} que le
résultat numérique. Voici ce qu'attend chaque consigne :
\end{propriete}

\begin{center}
\small
\begin{tabular}{ll}
\toprule
\textbf{Consigne} & \textbf{Ce qui est attendu} \\
\midrule
Calculer & Un résultat numérique, avec le détail du calcul qui y mène. \\
Démontrer / Montrer & Un raisonnement complet, citant le théorème ou la propriété utilisée. \\
Justifier & Une phrase précise expliquant \textbf{pourquoi} (pas seulement le résultat). \\
En déduire & Utiliser un résultat \textbf{déjà obtenu} sans le redémontrer depuis le début. \\
Vrai ou Faux & Répondre \textbf{et} justifier brièvement (sauf mention contraire). \\
Résoudre & Trouver toutes les solutions, avec la méthode qui y conduit. \\
Simplifier & Réduire une expression à sa forme la plus simple, en montrant les étapes. \\
Construire / Tracer & Un dessin \textbf{soigné}, réalisé aux instruments (règle, compas). \\
Déterminer / Donner & Donner la réponse demandée, avec le calcul si elle en découle. \\
Vérifier & Contrôler qu'une valeur donnée satisfait bien une condition, en le prouvant. \\
\bottomrule
\end{tabular}
\end{center}

\begin{remarque}
« Calculer » et « Déterminer » attendent presque toujours la \textbf{méthode}, pas
seulement la réponse finale : une réponse juste sans calcul visible perd des points, même
si le résultat est correct.
\end{remarque}

\section{Gérer son temps (2 heures)}

\begin{propriete}[Répartition indicative]
\begin{itemize}
\item \textbf{5 min} : lecture complète du sujet, repérage des parties les plus faciles.
\item \textbf{50 min} : Partie A (activités numériques puis géométriques), dans l'ordre ou
en commençant par ce qui est le plus maîtrisé.
\item \textbf{55 min} : Partie B (situation concrète), en lisant \textbf{deux fois}
l'énoncé avant de commencer (les données utiles sont souvent réparties sur plusieurs
paragraphes).
\item \textbf{10 min} : relecture, vérification des unités, des arrondis demandés, et de la
présentation.
\end{itemize}
\end{propriete}

\begin{remarque}
Sur une question bloquante, ne pas s'acharner plus de $3$ à $4$ minutes : passer à la
suite, et revenir dessus s'il reste du temps. Une question non traitée fait perdre ses
points ; s'obstiner dessus fait perdre aussi ceux des questions suivantes, non lues.
\end{remarque}

\section{Rédiger et présenter sa copie}

\begin{propriete}[Bonnes pratiques]
\begin{itemize}
\item \textbf{Encadrer ou souligner} chaque résultat final, avec son \textbf{unité}.
\item Recopier les \textbf{données de l'énoncé} avant de calculer (« On a $AB=5$ cm... »),
plutôt que d'écrire une formule sans expliquer d'où viennent les nombres.
\item Ne jamais laisser une question \textbf{blanche} : même une amorce de méthode
(formule correcte, début de calcul) peut rapporter un point.
\item Respecter \textbf{scrupuleusement} les arrondis et unités demandés (« arrondi au
centième », « en m$^{3}$ ») : une bonne valeur mal arrondie ou dans la mauvaise unité perd
des points.
\item Dans un problème concret (Partie B), \textbf{revenir au contexte} : conclure par une
phrase (« Le budget total est donc de... FCFA »), pas seulement un nombre isolé.
\end{itemize}
\end{propriete}

\section{Dix pièges fréquents à éviter}

\begin{synthese}[Les erreurs qui coûtent le plus de points]
\textbf{1.} $\dfrac12+\dfrac13\ne\dfrac25$ : additionner des fractions sans les réduire
d'abord au même dénominateur (chapitre 1).\\[3pt]
\textbf{2.} $a^{m}\times a^{n}=a^{m+n}$ (on \textbf{additionne}) $\ne(a^{m})^{n}=a^{m\times n}$
(on \textbf{multiplie}) (chapitre 2).\\[3pt]
\textbf{3.} $\sqrt{a+b}\ne\sqrt a+\sqrt b$, et $\sqrt{a^{2}}$ vaut $|a|$, pas $a$
(chapitre 3).\\[3pt]
\textbf{4.} $(a+b)^{2}\ne a^{2}+b^{2}$ : le double produit $2ab$ est oublié
(chapitre 4).\\[3pt]
\textbf{5.} Diviser ou multiplier une inéquation par un nombre \textbf{négatif} sans
inverser le sens de l'inégalité (chapitre 5).\\[3pt]
\textbf{6.} N'utiliser le théorème de Thalès que si le parallélisme est \textbf{donné} ou
déjà \textbf{prouvé} ; pour le prouver, c'est la réciproque qu'il faut utiliser (chapitre
10).\\[3pt]
\textbf{7.} Utiliser le théorème de Pythagore sur un triangle qui n'est \textbf{pas encore
prouvé} rectangle (chapitre 11).\\[3pt]
\textbf{8.} Confondre angle au centre et angle inscrit dans la relation « double » ; vérifier
qu'ils interceptent bien le même arc (chapitre 14).\\[3pt]
\textbf{9.} Confondre aire \textbf{latérale} et aire \textbf{totale} d'un solide (chapitre
15).\\[3pt]
\textbf{10.} Chercher la médiane sur une série \textbf{non ordonnée} (chapitre 16).
\end{synthese}

\begin{remarque}
Chacun de ces pièges est repris, avec un exemple corrigé, dans l'encadré « Piège
fréquent » de la fiche de révision du chapitre concerné.
\end{remarque}
